\section{Benchmark et avancement scientifique}
\label{sec:benchmark}

Le benchmark n'est pas une liste de scénarios mais un registre d'hypothèses.
Cette partie décrit ce qui est réellement mesuré, ce qui bloque l'inférence et
ce qui peut être publié aujourd'hui. Chaque chiffre est recalculé depuis un
artefact du dépôt, ou signalé comme lu dans un document.

\subsection{Le registre d'hypothèses et la porte V0}

Le registre place avant toute hypothèse une \emph{porte de validité} V0 : si une
tâche de contrôle soluble sans inférence échoue à l'un des deux bras, aucun
résultat en aval n'est lisible. Le seuil est gelé dans
\code{benchmarks/experiments/v0-control-manifest.json}, qui porte une valeur de
0,9, 24 étapes, 3 répétitions et la mention \code{saturating: true}. Le passage
de calibration rend un contrôle valide aux deux bras, trois fois sur trois. Le
registre en tire une porte \emph{partiellement franchie} : l'instrument est
lisible pour la parité d'exécution et pour un défaut répété, pas pour un défaut
rare.

Deux réserves doivent être écrites. D'une part, V0 ne couvre que MQTT : le
harnais SSH garde un dictionnaire d'identifiants fixe et n'est donc pas une
instance tenue à l'écart. D'autre part, et c'est un défaut de méthode, le seuil
n'est pas appliqué par la porte : la fonction de porte lit les quatre issues sur
« tous les essais réussis », tandis que le contrôle de seuil est calculé
séparément et ne conditionne rien. Avec trois répétitions, deux succès sur trois
valent 0,67, donc les deux formulations coïncident --- mais par coïncidence
numérique, non par application du seuil.

\begin{table}[H]
\centering
\caption{Statut du registre, recoupé avec les protocoles.}
\label{tab:registre}
\small
\begin{tabular}{@{}p{1.1cm}p{3.4cm}p{8.2cm}@{}}
\toprule
Item & Statut & Ce que cela signifie \\
\midrule
V0 & partiellement franchie & Contrôle valide sur MQTT, seuil gelé mais non appliqué par la porte ; contrôle saturant ; SSH non couverte \\
H1 & conditionnelle, non ouverte & La famille de chaînage d'identités n'existe pas dans le plan de couverture \\
H2 & mesurée le 22 octobre & Prédiction tenue, gain absorbé par un bras par convention \\
H3 & non implémentée & Les familles de contrat public sont planifiées \\
H4 & non ouverte & La famille de registres constructeur n'existe pas dans le corpus \\
H5 & conditionnelle, planifiée & Contre-épreuve de budget à écrire \\
H6 & prévue en sous-test & Greffable sur H2 ou H3, sans simulateur propre \\
C2 & rejetée & Deux relectures adverses, mort sur papier, aucun run \\
\bottomrule
\end{tabular}
\end{table}

\subsection{Le corpus : une strate logique minuscule}

Le catalogue énumère 29 scénarios. Le fichier dérivé
\code{benchmarks/labels/derived_index.json} compte 288 entrées attendues, dont
282 classées par le champ de type comme détectables par signature et 6 comme
relevant de la logique, soit 2,08~\%. Je les ai recomptés entrée par entrée.

Deux comptages coexistent et ne mesurent pas la même chose. La lecture
\emph{mécanique} range toutes les catégories sauf deux en signature. La lecture
\emph{interprétative} juge catégorie par catégorie si une règle bien conçue peut
conclure seule ; en reprenant ses catégories de signature sur l'index, on obtient
243 entrées sur 288, soit 84,4~\%, et les quatre lectures totalisent bien 288.
Les deux chiffres sont légitimes mais ne doivent jamais être cités l'un pour
l'autre.

\begin{table}[H]
\centering
\caption{Répartition des 288 entrées par catégorie et lecture interprétative.}
\label{tab:corpus}
\small
\begin{tabular}{@{}lrl@{}}
\toprule
Catégorie & Nombre & Lecture interprétative \\
\midrule
\code{misconfiguration} & 70 & signature \\
\code{data_exposure} & 54 & signature \\
\code{no_authentication} & 50 & signature \\
\code{default_credentials} & 41 & signature \\
\code{info_disclosure} & 16 & signature \\
\code{code_injection} & 13 & mixte \\
\code{insecure_update} & 12 & mixte \\
\code{weak_crypto} & 10 & mixte \\
\code{cve} & 9 & signature \\
\code{privilege_escalation} & 4 & chaînage \\
\code{broken_access_control} & 4 & logique \\
\code{missing_header} & 3 & signature \\
\code{auth_bypass} & 2 & logique \\
\midrule
Total & 288 & 243 signature, 35 mixte, 4 chaînage, 6 logique \\
\bottomrule
\end{tabular}
\end{table}

La strate logique n'occupe que trois scénarios, et 55 entrées seulement sont
rattachées à un chaînage d'au moins deux sauts. Un duel global sur ce corpus
mesurerait donc surtout la qualité des règles. C'est exactement ce que H2 a
observé.

\subsection{Les acquis réels}

\paragraph{H2, mesurée le 22 octobre.} Le corpus contient huit dossiers, deux par
groupe, pour 25 messages au total. Le run documenté exécute trois bras par
dossier, soit 24 bras par répétition et 72 bras pour les trois répétitions.
Sur les variantes brutes, le bras LLM obtient 0,5 --- un dossier gagné sur deux
---, le parseur déclaré 0,0, et le troisième bras, qui lit par convention sans
aucune clé du corpus, couvre 1,0. La prédiction tient au sens littéral, mais le
gain du modèle est absorbé par la convention.

\paragraph{C2, rejetée.} La revue conclut à l'arrêt : le premier candidat se
contredit, le second est mal rattaché et n'a que des issues non conformes, et le
comptage de familles n'est pas une mesure d'indépendance. Aucune fixture, aucun
run : c'est un résultat négatif de méthode, publiable comme limite.

\subsection{Le blocage de l'inférence}

L'outil d'inférence
(\code{src/benchmark/inference.py}) porte un défaut structurel. Les écarts valent
la différence entre les deux politiques ; or la p-value calculée est la
probabilité que la moyenne rééchantillonnée change de signe, non la p-value de
l'écart contre son seuil.

Je l'ai vérifié en exécutant la fonction : une série presque constante donne une
p-value nulle ; une série très dispersée de même moyenne donne aussi zéro ; une
série entièrement négative donne encore zéro. Quatre séries de comportements
différents produisent donc le même chiffre. En outre, le seuil n'entre dans
aucune décision : il ne sert qu'à produire un champ informatif. Enfin, en
dessous de trois familles, l'intervalle retourné est structurellement dégénéré.

\textbf{Aucune correction n'a eu lieu.} L'historique ne compte que deux commits
sur ce fichier, et toutes les branches disponibles portent le même calcul. Le
défaut est déjà écrit dans l'état courant du projet : l'outil produit
l'apparence d'une inférence qu'il ne fournit pas.

À ce défaut s'ajoute une impossibilité structurelle : la configuration d'un
essai inclut le commit, les empreintes et les budgets, donc deux familles
distinctes forment deux configurations distinctes. Chaque configuration ne
contient qu'une famille, et l'inférence en exige deux au minimum.

\subsection{L'admission de campagne}

Le module d'admission combine conditions de paire et conditions de campagne. La
condition d'adoption explicite des métriques de confirmation est un paramètre à
valeur par défaut fausse, jamais mis à vrai par le résumé ni par la ligne de
commande. Le module justifie ce verrou par un argument de contrat : les harnais
pilotes n'émettent pas ces compteurs. C'est exact de l'état actuel, mais ce n'est
pas une preuve d'insatisfiabilité : la condition redeviendrait vraie le jour où
un harnais publie ces compteurs. Aujourd'hui, elle n'est donc pas satisfaite.

Un second défaut, relevé pendant la rédaction puis vérifié, mérite relecture. La
condition de parité de budget exige que la première valeur existe et que les
deux soient égales ; l'absence de budget est donc correctement refusée, et ce
point est conforme à la règle annoncée. En revanche, le même module ne vérifie
pas que les budgets manquants sont distingués des budgets non comparables dans
les autres conditions ; la revue recommandait ce contrôle sans le retrouver.

\subsection{Ce qui est publiable aujourd'hui}

\textbf{Le publiable est la méthode, pas un résultat de supériorité.} Le dépôt ne
supporte aucune revendication de performance comparative : il n'existe ni
démonstration qu'un modèle bat un script compétent, ni intervalle de confiance
calculable, ni famille indépendante. Ce qui tient, en revanche, est consistant
et vérifiable :

\begin{itemize}
\item un registre d'hypothèses falsifiable, avec contraste déclaré avant la
fixture, oracle privé et mesure recalculée depuis le journal ;
\item une leçon de mesure sur H2 : une prédiction peut tenir au sens littéral et
ne rien démontrer, parce qu'un bras par convention absorbe le gain ;
\item un résultat négatif propre sur C2, documenté avant toute zone d'ombre ;
\item un inventaire honnête des trous : strate logique réduite, protocole
inconnu absent, bruit non encodé.
\end{itemize}

\subsection{Ce qui n'est pas démontré}

Ne doivent pas être écrits comme acquis : la supériorité du LLM ; un écart
significatif, puisque l'outil d'inférence produit une p-value sans rapport avec
l'écart et que publier un tel chiffre serait une faute ; la généralité, H2 ne
couvrant qu'un schéma, une question et un modèle ; le gain de coût, le pilote
exploratoire ne comptant qu'une paire sans intervalle ; la validité du seuil,
non appliqué par la porte ; et l'intégrité du corpus, le contrôle sain officiel
n'étant pas dans le catalogue.

\subsection{Les autres études actives}

Le plan V2 porte 44 fiches, 25 ensembles et plusieurs centaines de bras
conditionnels. Les qualifications réellement exécutées --- SSH en octobre, MQTT
quelques jours plus tard --- sont consignées comme qualifiant l'instrument, pas
un avantage. Le pilote à trois répétitions donne douze paires admissibles et ne
calcule aucun intervalle.

L'étude historique sur l'agent contre l'automatisation reste active et porte la
paire S1 et le programme sans modèle. Sa synthèse le dit sans détour : un
résultat descriptif réel, aucune supériorité générale établie. C'est la
formulation à conserver.
